% Körper – bank/koerper/ (zusammenbau v0.4, Heft)
\einheitenkopf[t2]{Körper}
\setcounter{aufgabe}{8}

% Anlauf (ohne Prüfkennung)
\begin{aufgabe}{Körper und Netze – Anlauf}
\begin{teile}
\teil Eine Konservendose hat diese Form. Welcher Körper ist es? \\ \kreuz{Kegel} \\ \kreuz{Zylinder} \\ \kreuz{Prisma}

\zylinder{1}{2.5}{}{}
\teil Ein Schuhkarton hat diese Form. Wie heißt der Körper?

\quader{4}{2}{2.5}{}{}{}
\leerfeld
\teil Ein Karton hat diese Form. Wie viele Ecken, Kanten und Flächen hat er?

\quader{4}{2}{2.5}{}{}{}
Ecken \leerfeld , Kanten \leerfeld , Flächen \leerfeld
\end{teile}
\end{aufgabe}

\begin{aufgabe}{Körper und Netze – Prüfungsaufgaben}
\begin{teile}
\teil Ein Zelt hat diese Form. Kreuze an, welcher Körper es ist \hfill (P10 2017 OS). \\ \kreuz{Pyramide} \\ \kreuz{Prisma} \\ \kreuz{Quader}

\prismadreieck{4}{2}{2.5}{}{}{}
\teil Ein Stück Käse hat diese Form. Kreuze an, welcher Körper es ist \hfill (P10 2017 OS). \\ \kreuz{Quader} \\ \kreuz{Pyramide} \\ \kreuz{Prisma}

\prismadreieck{3}{3}{1.5}{}{}{}
\teil Aus diesem Netz wird ein Körper gefaltet. Kreuze an, welcher Körper entsteht \hfill (P10 2018 OS). \\ \kreuz{Prisma} \\ \kreuz{Pyramide} \\ \kreuz{Quader}

\netzpyramide{2.5}{2}{}{}
\teil Ein Glasdach hat die Form einer Pyramide mit quadratischer Grundfläche, $2{,}4$ m hoch, Grundkante $3{,}6$ m. Zeichne die Höhe in das Schrägbild ein und beschrifte die Skizze mit beiden Maßen \hfill (P10 2015 OS).

\pyramide{3.5}{3.5}{3}{}{}{}
\teil Ein Zeltdach hat die Form einer Pyramide mit quadratischer Grundfläche, $1{,}9$ m hoch, Grundkante $2{,}8$ m. Zeichne die Höhe in das Schrägbild ein und beschrifte die Skizze mit beiden Maßen \hfill (P10 2015 OS).

\pyramide{3.5}{3.5}{2.5}{}{}{}
\teil Ein Schmuckstein hat die Form einer Pyramide mit quadratischer Grundfläche, $7$ cm hoch, Grundkante $9$ cm. Zeichne die Höhe in das Schrägbild ein und beschrifte die Skizze mit beiden Maßen \hfill (P10 2015 OS).

\pyramide{3.5}{3.5}{3.2}{}{}{}
\end{teile}
\end{aufgabe}

\begin{aufgabe}{Körper und Netze – Prüfungsaufgaben – weiter}
\begin{teile}
\teil Aus dem Netz wird eine Geschenkschachtel in Würfelform gefaltet. Die Fläche D ist der Deckel. Gib an, welche Fläche der Boden ist \hfill (P10 2016 OS).

$\begin{array}{cccc}  &  & A &  \\ D & N & M &  \\  &  & B & C \end{array}$
\leerfeld
\teil Das Netz wird zu einem Würfel gefaltet, der Würfel wird auf die Fläche D gestellt. Gib an, welche Fläche dann oben liegt \hfill (P10 2016 OS).

$\begin{array}{cccc}  & A &  &  \\  & M & N & D \\ C & B &  &  \end{array}$
\leerfeld
\teil Ein Pavillondach hat die Form einer Pyramide mit sechseckiger Grundfläche. Gib an, wie viele Kanten diese Pyramide hat \hfill (P10 2019 OS). \leerfeld Kanten
\teil Ein Briefbeschwerer hat die Form einer Pyramide mit dreieckiger Grundfläche. Gib an, wie viele Kanten diese Pyramide hat \hfill (P10 2019 OS). \leerfeld Kanten
\end{teile}
\end{aufgabe}

\begin{aufgabe}{Körper in ein Schrägbild einzeichnen (Kegel im Quader) – Prüfungsaufgaben}
\begin{teile}
\teil Ein Kegel aus Beton (Radius $25$ cm, Höhe $70$ cm) wird in einem quaderförmigen Karton verschickt, der so klein wie möglich ist. Skizziere den Kegel in das Schrägbild des Kartons \hfill (P10 2024 OS).

\quader{2.5}{2.5}{3.5}{}{}{}
\teil Eine Schultüte (Radius $18$ cm, Höhe $50$ cm) wird in einem quaderförmigen Karton verschickt, der so klein wie möglich ist. Skizziere den Kegel in das Schrägbild des Kartons \hfill (P10 2024 OS).

\quader{2.2}{2.2}{3}{}{}{}
\teil Ein Partyhut (Radius $9$ cm, Höhe $22$ cm) wird in einem quaderförmigen Karton verschickt, der so klein wie möglich ist. Skizziere den Kegel in das Schrägbild des Kartons \hfill (P10 2024 OS).

\quader{2.5}{2.5}{3}{}{}{}
\end{teile}
\end{aufgabe}

% Anlauf (ohne Prüfkennung)
\begin{aufgabe}{Quader – Anlauf}
\begin{teile}
\teil Wie viel Wasser passt in das Aquarium? Gefragt ist das … \kreuz{Volumen} \kreuz{Oberfläche}
\teil Quader: Länge $6$ cm, Breite $3$ cm, Höhe $2$ cm – Volumen? \leerfeld[cm³]
\teil Ein Würfel hat $5$ cm Kantenlänge – Volumen? \leerfeld[cm³]
\end{teile}
\end{aufgabe}

\begin{aufgabe}{Quader – Prüfungsaufgaben}
\begin{teile}
\teil Ein Würfel hat die Kantenlänge $8$ cm. Gib sein Volumen an \hfill (P10 2026 FOR). \leerfeld[cm³]
\teil Ein würfelförmiger Blumenkübel ist innen $2$ dm lang, breit und hoch. Gib an, wie viel Erde hineinpasst \hfill (P10 2026 FOR). \leerfeld dm³
\end{teile}
\end{aufgabe}

% Anlauf (ohne Prüfkennung)
\begin{aufgabe}{Prisma – Anlauf}
\begin{teile}
\teil Ein Zelt hat die Form dieses Prismas. Färbe die Grundfläche und zeichne die Höhe des Prismas farbig nach.

\prismadreieck{4}{2.5}{2}{}{}{}
\teil Prisma: Grundfläche $18$ cm², Höhe $7$ cm – Volumen? \leerfeld[cm³]
\teil Prisma mit rechteckiger Grundfläche $7$ cm × $3$ cm, Höhe $5$ cm – Volumen? Grundfläche \leerfeld[cm²] , Volumen \leerfeld[cm³]
\end{teile}
\end{aufgabe}

\begin{aufgabe}{Prisma – Prüfungsaufgaben}
\begin{teile}
\teil Ein Modell eines Deichs hat die Form eines Prismas. Die Grundfläche ist ein gleichschenkliges Trapez mit dem Flächeninhalt $2\,340$ cm², die Höhe des Prismas ist $85$ cm. Berechne das Volumen \hfill (P10 2014 OS). \leerfeld[cm³]
\teil Ein Blumenkasten hat die Form eines Prismas. Die Grundfläche ist ein gleichschenkliges Trapez mit dem Flächeninhalt $418$ cm², die Höhe des Prismas ist $72$ cm. Berechne das Volumen \hfill (P10 2014 OS). \leerfeld[cm³]
\teil Ein Brunnentrog hat die Form eines Prismas; Grund- und Deckfläche sind gleichseitige Dreiecke mit der Seite $1{,}80$ m und der Höhe $1{,}56$ m. Das Prisma ist $0{,}75$ m hoch. Berechne den Flächeninhalt der Deckfläche und weise nach, dass das Volumen etwa $1{,}1$ m³ beträgt \hfill (P10 2019 OS). Deckfläche \leerfeld[m²] , Volumen \leerfeld[m³]
\teil Ein Blumenkasten aus Beton hat die Form eines Prismas; Grund- und Deckfläche sind gleichseitige Dreiecke mit der Seite $0{,}70$ m und der Höhe $0{,}61$ m. Das Prisma ist $2{,}40$ m hoch. Berechne den Flächeninhalt der Deckfläche und weise nach, dass das Volumen etwa $0{,}5$ m³ beträgt \hfill (P10 2019 OS). Deckfläche \leerfeld[m²] , Volumen \leerfeld[m³]
\end{teile}
\end{aufgabe}

% Anlauf (ohne Prüfkennung)
\begin{aufgabe}{Zylinder – Anlauf}
\begin{teile}
\teil Eine Tasse ist oben $8$ cm breit, von Rand zu Rand durch die Mitte gemessen. Ist das Radius oder Durchmesser? Gib den Radius an. Radius \leerfeld[cm]
\teil Zylinder: Radius $5$ cm, Höhe $8$ cm – Volumen auf eine Stelle? \leerfeld[cm³]
\teil Zylinder: Durchmesser $8$ cm, Höhe $5$ cm – Volumen auf eine Stelle? \leerfeld[cm³]
\end{teile}
\end{aufgabe}

\begin{aufgabe}{Zylinder – Prüfungsaufgaben}
\begin{teile}
\steil Ein größerer Becher soll bei $9$ cm Höhe $600$ cm³ fassen. Berechne den Radius auf eine Stelle \hfill (P10 2022 OS). \leerfeld[cm]
\steil Ein Eimer soll bei $15$ cm Höhe $2\,000$ cm³ fassen. Berechne den Radius auf eine Stelle \hfill (P10 2022 OS). \leerfeld[cm]
\steil Eine Vase soll bei $18$ cm Höhe $1\,500$ cm³ fassen. Berechne den Radius auf eine Stelle \hfill (P10 2022 OS). \leerfeld[cm]
\teil Eine zylinderförmige Dose hat innen den Radius $3{,}5$ cm und die Höhe $9$ cm. Der Hersteller gibt „ca. 350 ml“ an. Weise nach, dass die Angabe stimmt \hfill (P10 2023 OS). Hinweis: $1$ cm³ $= 1$ ml.
\teil Ein Trinkbecher ohne Deckel ist innen $8{,}5$ cm hoch und hat den Radius $3{,}6$ cm. Prüfe, ob $350$ ml hineinpassen \hfill (P10 2022 OS). Hinweis: $1$ cm³ $= 1$ ml.
\teil Eine Konservendose hat innen den Radius $5{,}2$ cm und die Höhe $11{,}5$ cm. Auf dem Etikett steht „ca. 1 Liter“. Weise nach, dass die Angabe stimmt \hfill (P10 2023 OS).
\steil Eine Farbdose soll $2{,}5$ Liter fassen; ihr Deckel hat den Radius $7$ cm. Berechne die Innenhöhe der Dose auf eine Stelle \hfill (P10 2023 OS). \leerfeld[cm]
\steil Ein zylinderförmiger Messbecher fasst $0{,}75$ l; sein Boden hat den Radius $5$ cm. Berechne die Innenhöhe des Bechers auf eine Stelle \hfill (P10 2023 OS). \leerfeld[cm]
\end{teile}
\end{aufgabe}

% Anlauf (ohne Prüfkennung)
\begin{aufgabe}{Zusammengesetzt – Anlauf}
\begin{teile}
\teil Ein Haus mit Satteldach: Aus welchen zwei Körpern besteht es?
\teil Eine Betonstufe besteht aus zwei Quadern: unten $8$ dm × $4$ dm × $2$ dm, oben darauf $8$ dm × $2$ dm × $2$ dm. Volumen? \leerfeld dm³
\teil Ein Pool: Rechteck $7$ m × $4$ m, an der $4$-m-Seite ein Halbkreis mit $2$ m Radius, überall $1{,}6$ m tief. Welcher Term gibt das Wasservolumen in m³? \\ \kreuz{$7 \cdot 4 + \tfrac{1}{2} \cdot \pi \cdot 2^2 \cdot 1{,}6$} \\ \kreuz{$(7 \cdot 4 + \tfrac{1}{2} \cdot \pi \cdot 2^2) \cdot 1{,}6$} \\ \kreuz{$(7 \cdot 4 + \pi \cdot 2^2) \cdot 1{,}6$}
\end{teile}
\end{aufgabe}

\begin{aufgabe}{Zusammengesetzt – Prüfungsaufgaben}
\begin{teile}
\teil Ein zylinderförmiger Pflanzkübel (Radius $18$ cm, Höhe $55$ cm) soll in einem quaderförmigen Karton verschickt werden, der so klein wie möglich ist. Kreuze an und begründe \hfill (P10 2024 OS). \\ \kreuz{$19$ cm × $19$ cm × $56$ cm} \\ \kreuz{$37$ cm × $37$ cm × $56$ cm} \\ \kreuz{$37$ cm × $37$ cm × $53$ cm} \\ \kreuz{$56$ cm × $56$ cm × $56$ cm}
\teil Eine runde Torte (Radius $12{,}5$ cm, Höhe $9$ cm) soll in einem quaderförmigen Karton verschickt werden, der so klein wie möglich ist. Kreuze an und begründe \hfill (P10 2024 OS). \\ \kreuz{$13$ cm × $13$ cm × $11$ cm} \\ \kreuz{$26$ cm × $26$ cm × $8$ cm} \\ \kreuz{$26$ cm × $26$ cm × $11$ cm} \\ \kreuz{$32$ cm × $32$ cm × $11$ cm}
\teil Ein kegelförmiger Leitkegel (Radius $20$ cm, Höhe $70$ cm) soll in einem quaderförmigen Karton verschickt werden, der so klein wie möglich ist. Kreuze an und begründe \hfill (P10 2024 OS). \\ \kreuz{$41$ cm × $41$ cm × $71$ cm} \\ \kreuz{$21$ cm × $21$ cm × $71$ cm} \\ \kreuz{$41$ cm × $41$ cm × $69$ cm} \\ \kreuz{$71$ cm × $71$ cm × $71$ cm}
\steil Ein kegelförmiger Lampenschirm (Radius $18$ cm, Höhe $35$ cm) liegt in einer zylinderförmigen Verpackung (innen Radius $18{,}5$ cm, Höhe $36$ cm). Für den Kegel gilt $V = \tfrac{1}{3} \cdot \pi \cdot r^2 \cdot h$ ($r$ Radius, $h$ Höhe; Formelsammlung). Berechne das Restvolumen der Verpackung in cm³ \hfill (P10 2024 OS). \leerfeld[cm³]
\steil Eine Eistüte aus Schokolade ist ein Kegel (Radius $4$ cm, Höhe $11$ cm). Sie steckt in einer zylinderförmigen Hülle (innen Radius $4{,}2$ cm, Höhe $11{,}5$ cm). Für den Kegel gilt $V = \tfrac{1}{3} \cdot \pi \cdot r^2 \cdot h$ ($r$ Radius, $h$ Höhe; Formelsammlung). Berechne das Restvolumen der Hülle in cm³ auf eine Stelle \hfill (P10 2024 OS). \leerfeld[cm³]
\end{teile}
\end{aufgabe}
